%% kapa: exam=1 round=22 year=2011 period=2 form=A subject=accounting subject_ko=회계학 seq=31 number=31
%% source: Q-Net 감정평가사 기출문제 (cst00309) artlSeq=5036790 file=extracted/1008/5036790/2011년 제22회 감정평가사 1차시험 기출문제/2교시.hwp
%% conversion: kordoc 4.24.51 -> tools/gen1.py (auto); review=proofread-sonnet; flags=none
\begin{question}{31}
(주)서울은 두 종류의 제품(컴퓨터와 프린터)을 생산하고 있다. 회사의 제조활동은 다음 4가지로 구분되며, 활동별 제조간접원가와 관련된 자료는 다음과 같다.

{\small\centering
\begin{tabular}{|c|c|c|r|}
\hline
활동 & 원가동인 & \makecell{연간\\원가동인 수} & \multicolumn{1}{c|}{연간 활동별 제조간접원가} \\
\hline
\makecell{생산\\준비} & 생산준비시간 & 600시간 & ￦900,000 \\
\hline
\makecell{재료\\이동} & 재료이동횟수 & 1,800회 & 1,080,000 \\
\hline
\makecell{기계\\사용} & 기계작업시간 & 400시간 & 1,200,000 \\
\hline
\makecell{수선\\유지} & 기계작업시간 & 400시간 & 800,000 \\
\hline
\end{tabular}\par}

컴퓨터에 대한 생산량 및 원가자료가 다음과 같을 때, 활동기준원가계산(ABC)에 의한 컴퓨터의 단위당 제조원가는? \par
생산량2,000단위 \par
생산준비시간 300시간 \par
재료이동횟수 1,170회 \par
기계작업시간 250시간 \par
단위당 직접재료원가 ￦3,000 \par
단위당 직접노무원가 ￦4,000 \par
\begin{choices}
  \item ￦7,562
  \item ￦8,201
  \item ￦8,932
  \item ￦9,653
  \item ￦10,052
\end{choices}
\end{question}
\answer{2}
